\documentclass[10pt,a4paper]{amsart}
\usepackage[margin=19mm]{geometry}
\usepackage{amsmath,amssymb,mathtools}
\usepackage[scr=rsfs,cal=euler]{mathalfa}
\usepackage{xfrac}
\usepackage[unicode,hidelinks,bookmarks=false]{hyperref}
\newcommand{\ie}{i.e.\ }
\newcommand{\cf}{cf.\ }
\newcommand{\resp}{respectively\ }
\newcommand{\Subsection}{\subsection}
\providecommand{\nobreakdash}{\nobreak\hbox{-}}
\setlength{\emergencystretch}{2em}
\multlinegap0pt
\usepackage{amssymb}
\usepackage{xcolor}
\usepackage{tikz}
\usepackage{mathtools}
\usepackage{float}
\newtheorem{defn}{Definition}
\newtheorem{lem}{Lemma}
\newcommand{\Liq}{\mathbb L}
\newcommand{\Q}{\mathbb Q}
\newcommand{\cC}{\mathcal{C}}
\newcommand{\cI}{\mathcal{I}}
\newcommand{\cJ}{\mathcal{J}}
\newcommand{\cN}{\mathcal{N}}
\renewcommand{\geq}{\geqslant}
\renewcommand{\leq}{\leqslant}
\title{Maximal universal invariants of knots from finite quotients of Verma modules}
\author{Cristina Ana-Maria Anghel, Jun Murakami}
\date{Source: arXiv:2606.19113v1 (CC BY 4.0)}
\begin{document}
\maketitle

\noindent\emph{Source and licence.} Maximal universal invariants of knots from finite quotients of Verma modules. Version arXiv:2606.19113v1 (CC BY 4.0), distributed by arXiv under the Creative Commons Attribution 4.0 International licence, http://creativecommons.org/licenses/by/4.0/, as recorded in the arXiv licence metadata for this exact version and shown on the abstract page https://arxiv.org/abs/2606.19113v1. Full licence notice: \texttt{licenses/arxiv-2606.19113-CC-BY-4.0.txt}.\par\smallskip

\noindent\textit{Editorial scope.} Counted unit: the article's lem lem:EQ3 (source0.tex lines 1427--1433). The article's complete original proof of it is delivered with it (source0.tex lines 1435--1615, one \texttt{proof} environment of 181 lines). No mathematics is written, repaired or re-derived: the statement and the proof are the article's own lines, in the article's own notation, and no retained line is substituted or re-flowed.  Retained uncounted as the source-local context the proof needs: lines 370--373; lines 728--748; lines 751--759; lines 1617--1620.  Nothing was omitted from any retained span, so the package carries no omission record.  No retained line contains a citation, so the wrapper carries no bibliography and no bibliography entry.  No retained line contains a figure environment, an image-inclusion command or drawing code, so no image is delivered and the package declares no asset dependency.  Editorial formatting only, in addition: an AMS \texttt{amsart} preamble that restates the article's own declarations and macros used by the retained lines, namely the environments \texttt{defn}, \texttt{lem} on separate counters, together with the standard packages those lines need.  The retained environments print in this document as Definition 1, Lemma 1 in order of appearance, whereas the article prints the same items with the numbering of its own full document.  The complete native source of the article as distributed in the arXiv e-print for this exact version is bundled unchanged as \texttt{sources/203117/source0.tex}.
\begin{equation}\label{qideal}
\cI_\cN:=\left(\varphi_\cN(q^2)\cdot (s^2q^{-2(\cN-1)}-1)\right)
\subseteq \Liq=\Q[q^{\pm2}, s^{\pm 2}]
\end{equation}

\begin{defn}[Condition ideal] 
\label{def:condition} 
Let us consider the following ideals in $\Liq$: 
%
%\begin{equation}
%\cI_{\cN}:=\Big( \varphi_\cN(q^2)\cdot (s^2q^{-2\cN+2}-1) \Big),
%\end{equation}
%
\begin{equation*}
\cJ_{\cN,k}:=\left( \varphi_k(q^2),\  \ \prod_{i=0}^{k-2}(s^2q^{-2i}-1) \right)
\end{equation*}
for $k| \cN$, $k \neq 1$, $2$, $\cN-1$, $\cN$ where $\varphi_k$ is the $k^{th}$ cyclotomic polynomial, and
\begin{equation*}
\widetilde{\cJ}_\cN = \cI_\cN \cap\!\!\!\! \bigcap_{\text{\scriptsize$\begin{matrix}
d | \cN\\[-1pt]
d\neq 1, \cN
\end{matrix}$}} \cJ_{\cN,d},  
\end{equation*}
where $\cI_\cN = \Big( \varphi_\cN(q^2)\cdot (s^2q^{-2\cN+2}-1) \Big)$ given by \eqref{qideal}.  
We call $\widetilde{\cJ}_\cN$ the ``condition ideal''.
\end{defn}

\begin{defn}[Maximal quotient ideal]
\label{def:maximal}
Let us consider as well the ideal
\begin{equation}\label{idealb}
\cC_\cN = \sum_{\substack{i,j=0 \\ j+n\geq \cN,0\leq  n\leq i}}^{\cN-1}
 \left( {n+j \brack j}_q \prod_{k=0}^{n-1}(sq^{-k-j}-s^{-1}q^{k+j}) \right) \subseteq \Liq_{\cN}
\end{equation}
and call it the ``maximal quotient ideal''.
\end{defn}

\begin{lem}\label{lem:qi}
If $k$ and $l$ are both multiple of $d$, then ${k \brack l}_q$
is equal to ${k/d \brack l/d}_1$ or $-{k/d \brack l/d}_1$ modulo $\Phi_d(q^2)$.  
\end{lem}

\begin{lem}
\label{lem:EQ3}
The   three  structural ideals introduced in Definitions \ref{def:condition} and \ref{def:maximal} are equal, i.e. 
\begin{equation*}
\cC_\cN = \widetilde{\cJ}_{\cN}. 
\end{equation*}
\end{lem}

\begin{proof}
We use the following intermediate ideal 
%
\begin{equation*}%\label{eq:INprime}
\cI'_\cN = \sum_{k=0}^{\cN-2} \left( {\cN \brack k+1}_q \cdot \prod_{i=1}^{\cN-k-1}(s^2q^{-2k-2i}-1) \right)\subseteq \Liq.
\end{equation*}
%
We first show that we have the equality $\cI_{\cN}^\prime = \widetilde{\cJ}_{\cN}$, in two parts, as below.
\par
{\bf Step I: $\cI_{\cN}^\prime \subset \widetilde{\cI}_{\cN}$} 
\par  
Each ideal $\left( {\cN \brack k+1}_q \cdot \prod_{i=1}^{\cN-k-1}(s^2q^{-2k-2i}-1) \right)$ in 
$\cI_\cN^\prime$ contains the factor $(s^2 q^{-2\cN+2}-1)$, so 
\begin{equation}\label{r1'}
\cI^\prime_\cN = (s^2 q^{-2\cN+2}-1)\left(\sum_{k=0}^{\cN-2}
 \Big({\cN \brack k+1}_q \cdot \prod_{i=1}^{\cN-k-2}(s^2q^{-2k-2i}-1) \Big)\right).
\end{equation} 
Now we look at the ideal $\sum_{k=0}^{\cN-2}
 \Big({\cN \brack k+1}_q \cdot \prod_{i=1}^{\cN-k-2}(s^2q^{-2k-2i}-1) \Big)$. 
 %
 \begin{equation}\label{r2'}
 \begin{aligned} \sum_{k=0}^{\cN-2}
 \Big({\cN \brack k+1}_q \cdot \prod_{i=1}^{\cN-k-2}(s^2q^{-2k-2i}-1) \Big)
 &=
\sum_{k=1}^{\cN-1}
  \Big({\cN \brack k}_q \cdot \prod_{i=1}^{\cN-k-1}(s^2q^{-2k-2i+2}-1) \Big)
 =\\
&=\sum_{k=1}^{\cN-1}
  \Big({\cN \brack k}_q \cdot \prod_{i=0}^{\cN-k-2}(s^2q^{-2k-2i}-1) \Big).
 \end{aligned}
 \end{equation}
 %
 Now, for the term associated to $k=\cN-2$ we write the quantum number $[\cN]_{q}$ as a product of all $\varphi_d(q^2)$ for $d\mid \cN, d \neq 1$, and so we obtain the following description of the ideals:
 \begin{equation*}
 \begin{aligned}
&\Big(\bigcap_{\text{\scriptsize$\begin{matrix}
 d\mid \cN \\ d \neq 1 \end{matrix}$}} \Big( \varphi_d(q^2)\Big)\Big) +
\sum_{k=1}^{\cN-2}
  \Big({\cN \brack k}_q \cdot \prod_{i=0}^{\cN-k-2}(s^2q^{-2k-2i}-1) \Big) 
  =\\
 &=\bigcap_{\text{\scriptsize$\begin{matrix}
 d\mid \cN \\ d \neq 1 \end{matrix}$}} \left(\Big( \varphi_d(q^2)\Big) +
\sum_{k=1}^{\cN-2}
  \Big({\cN \brack k}_q \cdot \prod_{i=k}^{\cN-2}(s^2q^{-2i}-1) \Big)  \right).
 \end{aligned}
 \end{equation*}
 The cyclotomic plynomial $\varphi_d(q^2)$ divide ${\cN \brack k}_q$ if and only if $k$ is not a multiple of $d$.  
 Moreover, $\varphi_d(q^2)$ and ${\cN \brack ld}_q$ are mutually prime.  
 We also know that $q^{2d}$ is equal to $1$ modulo $\varphi_d(q^2)$.  
 Therefore, we have
 %
\begin{equation*}
\Big( \varphi_d(q^2)\Big) +
\sum_{k=1}^{\cN-2}
  \Big({\cN \brack k}_q \cdot \prod_{i=k}^{\cN-2}(s^2q^{-2i}-1) \Big)  
=
\Big( \varphi_d(q^2)\Big) +
\sum_{l=1}^{\cN/d-1}
  \Big({\cN \brack ld}_q \cdot \prod_{i=ld}^{\cN-2}(s^2q^{-2i}-1) \Big). 
\end{equation*}
Replacing $l$ by $\cN/d - l$, we get
\begin{equation*}
\Big( \varphi_d(q^2)\Big) +
\sum_{l=1}^{\cN/d-1}
  \Big({\cN \brack ld}_q \cdot \prod_{i=ld}^{\cN-2}(s^2q^{-2i}-1) \Big)
  =
\Big( \varphi_d(q^2)\Big) +
\sum_{l=1}^{\cN/d-1}
  \Big({\cN \brack ld}_q\prod_{i=0}^{ld-2}(s^2q^{-2i}-1) \Big) ,
\end{equation*}
which is contained in $\left(\varphi_d(q^2), \  \ \prod_{i=0}^{d-2}(s^2q^{-2i}-1) \right)$. 
%
Hence, $\cI_\cN^\prime$ is contained in $\widetilde{\cJ}_\cN$.
%
\par
{\bf Step II: $\widetilde{\cJ}_\cN \subseteq \cI_\cN^\prime$}
\par
By Lemma \ref{lem:qi} given below, we have ${\cN \brack ld}_q = \pm{\cN/d \brack l}_1$, which is a non-zero integer and is invertible in $\mathbb{Q}$, so we obtain:
\begin{equation*}
\begin{aligned}
\Big( \varphi_d(q^2)\Big) +
\sum_{l=1}^{\cN/d-1}
  \Big({\cN \brack ld}_q\prod_{i=0}^{ld-2}(s^2q^{-2i}-1) \Big) 
    &=
 \Big( \varphi_d(q^2)\Big) +
\sum_{l=1}^{\cN/d-1}
  \Big(\prod_{i=0}^{ld-2}(s^2q^{-2i}-1) \Big) 
  =\\
  &=
   \Big( \varphi_d(q^2)\Big) +
  \Big(\prod_{i=0}^{d-2}(s^2q^{-2i}-1) \Big).
 \end{aligned}
\end{equation*} 
Then, following relations \eqref{r1'} and \eqref{r2'} we have:
\begin{equation*}
\begin{aligned}
\cI^\prime_\cN &= (s^2 q^{-2\cN+2}-1)\left(\sum_{k=0}^{\cN-2}
 \Big({\cN \brack k+1}_q \cdot \prod_{i=1}^{\cN-k-2}(s^2q^{-2k-2i}-1) \Big)\right)=\\
 &=(s^2 q^{-2\cN+2}-1) \cdot \bigcap_{\text{\scriptsize$\begin{matrix}
 d\mid \cN \\ d \neq 1 \end{matrix}$}} \Big(\Big( \varphi_d(q^2)\Big) +
  \Big(\prod_{i=0}^{d-2}(s^2q^{-2i}-1) \Big) \Big).
 \end{aligned}
\end{equation*} 
On the other hand, let us recall that 
\begin{equation*}
\widetilde{\cJ}_\cN =   
\left((s^2 q^{-2\cN+2} - 1) \cdot\left(\varphi_\cN(q^2)\right)\right) \cap
 \bigcap_{\text{\scriptsize$\begin{matrix}
 d\mid \cN \\ d \neq 1, \cN \end{matrix}$}}
 \left(\varphi_d(q^2),\ \  \prod_{i=0}^{d-2}(s^2q^{-2i}-1) \right)
\end{equation*}
%
The previous relations imply that $\widetilde{\cI}_\cN \subseteq \cI_\cN^\prime$ and so we conclude the equality $\cI_\cN^\prime = \widetilde{\cI}_\cN$.  
%
%
\par
Next, we will show that $\cI_{\cN}^\prime = \cC_{\cN}$ in two steps.
%
\par
{\bf Step III: $\cI_{\cN}^\prime \subset \cC_{\cN}$ } 
\par
For this, we start with the definition of $\cC_\cN$:
\begin{multline*}
\cC_\cN= \sum_{\substack{i,j=0 \\ j+n\geq \cN,0\leq  n\leq i}}^{\cN-1}
 \left( {n+j \brack j}_q \prod_{k=0}^{n-1}(sq^{-k-j}-s^{-1}q^{k+j}) \right) 
 \\
 =
 \sum_{i=0}^{\cN-1}
 \sum_{j=1}^{\cN-1}
 \sum_{n=\cN-j}^{i}
 \left( {n+j \brack j}_q \prod_{k=0}^{n-1}(sq^{-k-j}-s^{-1}q^{k+j}) \right) 
 \\=
 \sum_{i=0}^{\cN-1}
 \sum_{j=0}^{\cN-2}
 \sum_{n=\cN-j-1}^{i}
 \left( {n+j+1 \brack j+1}_q \prod_{k=0}^{n-1}(sq^{-k-j-1}-s^{-1}q^{k+j+1}) \right).  
\end{multline*}
For $n=\cN-j-1$, 
\begin{multline*}
 \left( {n+j+1 \brack j+1}_q \prod_{k=0}^{n-1}(sq^{-k-j-1}-s^{-1}q^{k+j+1}) \right)  
 =
 \\
  \left( {\cN \brack j+1}_q \prod_{k=0}^{\cN-j-2}(sq^{-k-j-1}-s^{-1}q^{k+j+1}) \right)  
  =
  \\
   \left( {\cN \brack j+1}_q \prod_{k=1}^{\cN-j-1}(sq^{-k-j}-s^{-1}q^{k+j}) \right).
\end{multline*}
Therefore, we obtain that $\cI_{\cN}^\prime \subset \cC_{\cN}$.  
\par
%
{\bf Step IV: $\cC_{\cN}\subset \cI_{\cN}^\prime$} 
%
\par
For this last step we are going to show that:
\begin{equation*}
\left( {n+j+1 \brack j+1}_q \prod_{k=0}^{n-1}(sq^{-k-j-1}-s^{-1}q^{k+j+1}) \right) 
\in \cI_\cN^\prime.
\end{equation*}
Let $l = n+j+1 - \cN$.  
Then
\begin{equation*}
\left( {n+j+1 \brack j+1}_q \prod_{k=0}^{n-1}(sq^{-k-j-1}-s^{-1}q^{k+j+1}) \right)
=
\left( {\cN+l \brack j+1}_q \prod_{k=1}^{\cN+l-j-1}(sq^{-k-j}-s^{-1}q^{k+j}) \right).
\end{equation*}
We show that the above is contained in $\cI_\cN^\prime$ by an induction on $l$.  
If $l=0$, then it is contained in $\cI_\cN^\prime$.  
Now we assume that $\left( {\cN+l \brack j+1}_q \prod_{k=1}^{\cN+l-j-1}(sq^{-k-j}-s^{-1}q^{k+j}) \right) \in \cI_\cN^\prime$.  
Then
\begin{multline*}
\left( {\cN+l+1 \brack j+1}_q \prod_{k=1}^{\cN+l-j}(sq^{-k-j}-s^{-1}q^{k+j}) \right)
=
\\
q^{\cN+l-j}\left( {\cN+l \brack j+1}_q \prod_{k=1}^{\cN+l-j}(sq^{-k-j}-s^{-1}q^{k+j}) \right)
+
\\
q^{-j-1}\left( {\cN+l \brack j}_q \prod_{k=1}^{\cN+l-j}(sq^{-k-j}-s^{-1}q^{k+j}) \right).  
\end{multline*}
The last  two terms are contained in $\cI_\cN^\prime$ and so
$\left( {\cN+l+1 \brack j+1}_q \prod_{k=1}^{\cN+l-j}(sq^{-k-j}-s^{-1}q^{k+j}) \right) \in \cI_\cN^\prime$.  
\end{proof}

\end{document}
